\chapter{Methodology}
\label{chap:methodology}

\section{Introduction and Scope}
This chapter sets forth a \emph{multi-phase methodology} aimed at refining automated ICD-10-CM code prediction, \emph{especially} for rare codes. A comprehensive framework is proposed, featuring an \emph{ontology-based disease grouping} strategy, a pipeline for \emph{synthetic discharge summary generation}, a \emph{hybrid model training} approach that integrates both real and synthetic data, and an evaluation scheme specifically \emph{attuned} to rare code distribution imbalances. A \emph{modular and extensible workflow} is emphasized, so that the methodology may be adapted readily to alternative scenarios involving imbalanced classification. Ultimately, validation is planned on a real-world clinical dataset to assess the \emph{generalizability} and \emph{utility} of this framework.

\section{Ontology-Based Disease Grouping}
Rare ICD-10-CM codes frequently exhibit \emph{extremely low prevalence}, which hinders the ability of predictive models to learn discriminative features. In order to mitigate this challenge, an \emph{ontology-based grouping} approach has been envisioned. This process entails parsing established medical ontologies---such as \emph{SNOMED CT}, \emph{UMLS}, and the \emph{Orphanet Rare Disease Ontology}---to merge \emph{hierarchical} and \emph{associative} relationships.

It is intended that each rare disease will be anchored in a \emph{knowledge graph}, and \emph{closely related} concepts will be identified via direct graph traversal. These items will be clustered into a \emph{unified disease family}, thereby providing a \emph{contextual} lens through which the rare disease can be understood. By associating each rare disease with more common conditions, the model will receive exposure to realistic clinical context, improving its capacity to recognize the \emph{linguistic cues} relevant to rare diagnoses. Additionally, ontology mapping ensures that no \emph{orphaned codes} are disregarded, given that each disease is matched at least to a more general ICD-10-CM code if an exact match proves elusive.

\emph{It is anticipated} that this ontology-centric approach will enable more \emph{robust learning} for rare codes, because they will be presented alongside clinically related diagnoses. By placing rare codes in \emph{medically coherent families}, the methodology reduces the isolation that often occurs when codes appear only a handful of times in conventional datasets.

\section{Synthetic Discharge Summary Generation}
Because rare codes remain \emph{infrequent} in standard patient records, a \emph{synthetic data generation} process is introduced to amplify the presence of these \emph{underrepresented labels}. A retrieval-augmented generation pipeline has been planned, wherein medical facts, professional guidelines, and relevant textual fragments are \emph{retrieved} to guide a large language model in producing draft discharge summaries.

\begin{itemize}
    \item \textbf{Structured Prompting:} Placeholders will be embedded for sections such as \emph{Chief Complaint}, \emph{History of Present Illness}, \emph{Hospital Course}, \emph{Medications}, and \emph{Discharge Diagnoses}.
    \item \textbf{Disease Clusters:} Rare diseases will be linked with related diagnoses identified through ontology-based grouping, ensuring that \emph{realistic co-occurrences} appear in the synthetic text.
    \item \textbf{Multi-Agent Feedback Loop:}
    \begin{itemize}
        \item A \emph{draft discharge summary} is generated by the language model.
        \item A \emph{critic agent} then examines the draft for logical inconsistencies and \emph{missing details}.
        \item Iterative refinements are made until the note meets \emph{plausibility} and \emph{consistency} criteria.
    \end{itemize}
\end{itemize}

The final \emph{synthetic discharge summaries} must provide sufficient evidence for the relevant rare ICD-10-CM codes, mirroring the style and structure of authentic clinical documentation. This balanced approach is expected to \emph{significantly} enrich the training dataset for rare codes, thereby reducing distributional skew and enhancing overall model recall for these infrequent conditions.

\section{Deep Learning Model Training and Testing}
An \emph{augmented training corpus} composed of both real and synthetic discharge summaries will subsequently be employed in training two deep learning architectures: (i) a \emph{transformer-based} model, and (ii) a \emph{CNN-based} model. The transformer-based model is \emph{anticipated} to capture \emph{long-range dependencies}, whereas the CNN-based model should excel at identifying \emph{localized textual patterns} pertinent to specific codes.

A multi-label classification scheme will be used, assigning probabilities to all potential ICD-10-CM codes. To address \emph{class imbalance}, strategies such as \emph{weighted sampling}, \emph{focal loss}, and \emph{class-weighted loss} have been planned. Synthetic summaries will be \emph{oversampled} or \emph{weighted} more heavily, ensuring that rare codes appear frequently enough in training batches. The training procedure will commence with \emph{binary cross-entropy} and moderate class weighting, followed by the incremental introduction of \emph{focal loss} during later epochs, thus allowing the model to \emph{fine-tune} its attention on the \emph{most difficult} labels.

A development set containing both real and synthetic notes will guide \emph{hyperparameter tuning} and \emph{early stopping}. To confirm the \emph{practical value} of this method, the final evaluation will be conducted on a \emph{held-out portion} of real MIMIC-IV discharge summaries, assessing whether improvements in \emph{rare code classification} carry over to genuine patient data.

\section{Evaluation Metrics and Analysis}
The planned evaluation seeks to capture \emph{various dimensions} of performance, with special focus on rare code detection. Several metrics will be reported, each \emph{illuminating} a different aspect of classifier efficacy. The following formulas present the principal metrics in detail (where $TP$ = true positives, $FP$ = false positives, $FN$ = false negatives, and $C$ = total number of classes):

\begin{enumerate}
    \item \textbf{Micro-F1}
    \begin{equation}
        \text{Micro-F1} = 2 \times 
        \frac{\sum_{i=1}^{C} TP_{i}}
        {\sum_{i=1}^{C} (2 \cdot TP_{i} + FP_{i} + FN_{i})}
    \end{equation}
    \emph{Micro-F1} aggregates the contributions of all codes to compute the overall precision and recall across the entire label space.

    \item \textbf{Macro-F1}
    \begin{equation}
        \text{Macro-F1} = \frac{1}{C} \sum_{i=1}^{C} \text{F1}_{i}
    \end{equation}
    Each class-level F1 score $\text{F1}_{i}$ is calculated, and their average is taken to highlight how the model performs on \emph{each individual} label, giving rare codes equal weight to common codes.

    \item \textbf{Recall@\emph{k}}
    \begin{equation}
        \text{Recall@k} = 
        \frac{\text{Number of times the correct code is in the top } k \text{ predictions}}
        {\text{Total number of samples}}
    \end{equation}
    This metric indicates the fraction of cases in which the \emph{true} codes appear among the top $k$ predictions, a setting relevant in \emph{coding-assistance systems} that present a limited list of suggestions.

    \item \textbf{Exact Match Ratio}
    \begin{equation}
        \text{Exact Match} = 
        \frac{\text{Number of notes for which the full set of codes is correctly predicted}}
        {\text{Total number of notes}}
    \end{equation}
    This stricter criterion gauges the extent to which all predicted codes align \emph{perfectly} with the ground truth set for each note.

    \item \textbf{Never-Predicted Rate}
    \begin{equation}
        \text{Never-Predicted Rate} = 
        \frac{\text{Number of codes that are never predicted by the model}}{C}
    \end{equation}
    This measure highlights whether any codes---particularly \emph{rare ones}---are completely ignored.
\end{enumerate}

\emph{Emphasis} will be placed on determining whether the \emph{ontology-based grouping} and \emph{synthetic data} augmentations produce measurable gains in macro-F1 and reduce the \emph{never-predicted rate}. A targeted error analysis will be conducted for codes that remain \emph{unrecalled}, providing directions for further refinements.

\section{Processing MIMIC-IV as Our Main Test Dataset}
\label{sec:processing_mimic_iv}

In this work, the MIMIC-IV database will be employed as the principal test dataset for automated ICD coding. This dataset contains de-identified clinical information collected from the Beth Israel Deaconess Medical Center (BIDMC), with particular emphasis on structured hospital records and unstructured discharge summaries. The following steps describe how the data will be prepared, managed, and validated in a manner consistent with privacy regulations and best practices in clinical data science.

\subsection{Data Extraction and Integration}
The relevant tables from MIMIC-IV (e.g., \texttt{admissions}, \texttt{patients}, \texttt{diagnoses\_icd}, and \texttt{discharge}) will be identified and retrieved. These tables will then be joined on keys such as \texttt{subject\_id} and \texttt{hadm\_id}, thus creating a unified view of patient admissions, annotated diagnoses, and the corresponding discharge narratives. During this process, all admissions lacking a corresponding discharge summary or valid ICD-10 code assignments will be flagged for potential exclusion, ensuring that the resulting dataset remains suitable for supervised learning.

\subsection{Text Standardization and Cleaning}
Each discharge summary will be examined to remove any remaining placeholders or noise introduced by the de-identification process. Lowercasing and tokenization will be conducted if deemed beneficial for subsequent analyses, and the presence of medical abbreviations or numeric identifiers will be noted. Segmentation into sections (such as “History of Present Illness” and “Medications”) may be carried out to reduce text complexity, though the entire narrative will remain intact if full-length modeling is desired.

\subsection{Handling Duplicates and Missing Data}
Instances in which multiple, nearly identical discharge summaries exist for a single admission will be managed by detecting duplicates through textual similarity checks. These checks will reduce the risk of overrepresenting particular patients or hospital stays. Any admissions without sufficiently long or clearly documented notes will be handled through either imputation strategies or careful removal, thereby minimizing noise in model training and evaluation.

\subsection{Code Distribution and Rare Labels}
A preliminary examination of the ICD-10 codes assigned to each admission will be performed to understand their frequency distribution. Since ICD codes often follow a long-tail pattern, additional attention will be devoted to codes that appear infrequently. If these rare labels contribute significantly to the overall error rate, strategies such as targeted oversampling or retrieval-augmented data enrichment may be considered, provided such techniques respect the constraints of de-identification and data plausibility.

\subsection{Database Management and Version Control}
All data transformations and table joins will be executed within a controlled environment (e.g., DuckDB) to allow for efficient querying and seamless integration with Python-based analytic workflows. The resulting datasets, along with any derived tables, will be tracked using version control to ensure that incremental changes are systematically logged and reproducible. This practice will enable transparent experimentation and simplify the process of rolling back to earlier dataset states if necessary.

\subsection{Security and Compliance}
All patient information in MIMIC-IV is protected via de-identification, and no direct identifiers (e.g., names, addresses) are present. It will be confirmed that the data use agreement from PhysioNet is upheld, particularly regarding re-identification risks and data sharing restrictions. Any synthetic expansions or textual outputs generated by the pipeline will be carefully reviewed to ensure that confidential details are not reintroduced.

